\documentclass[11pt]{article}
\usepackage[letterpaper,margin=1in]{geometry}
\usepackage{amsmath,amssymb,amsthm}
\usepackage{booktabs,tabularx,array}
\usepackage[colorlinks=true,linkcolor=black,citecolor=black,urlcolor=blue]{hyperref}
\hypersetup{
  pdftitle={A nine-vertex counterexample to Jensen-Toft Problem 5.8},
  pdfauthor={Alex Chan}}

\newtheorem{theorem}{Theorem}
\newtheorem{lemma}{Lemma}
\newcommand{\out}{N^{+}}

\title{A nine-vertex counterexample to\\Jensen--Toft Problem 5.8}
\author{Alex Chan\\[2pt]
  \small Harvard University and National Bureau of Economic Research\\[-1pt]
  \small \texttt{achan@hbs.edu}}
\date{\today}

\begin{document}
\maketitle

\begin{abstract}
We give a negative answer to the question in Jensen and Toft,
\emph{Graph Coloring Problems}, Section~5.8, asking whether every
noncomplete $k$-critical graph, for $k\geq4$, contains a noninduced
$(k-1)$-critical subgraph. The witness is the classical nine-vertex
5-Ore graph $O_5$. We show that its 4-critical subgraphs are exactly
three copies of $K_4$ and twelve copies of the Moser spindle, all
induced. The proof is elementary. A Lean formalization independently
certifies the counterexample property.
\end{abstract}

\section{The question}

All graphs are finite and simple. A graph is \emph{$k$-critical} if
its chromatic number is $k$ and every proper subgraph is
$(k-1)$-colorable. A subgraph $H$ of $G$ is \emph{induced} if
$H=G[V(H)]$.

Jensen and Toft~\cite[Section~5.8, p.~103]{JT} ask whether every
noncomplete $k$-critical graph, for $k\geq4$, contains a
$(k-1)$-critical subgraph that is not induced. We answer this
question as stated.

\begin{theorem}\label{thm:main}
There is a noncomplete 5-critical graph $G$ on nine vertices and
nineteen edges whose 4-critical subgraphs are exactly fifteen
induced subgraphs: three copies of $K_4$ and twelve copies of the
Moser spindle. In particular, $G$ has no noninduced 4-critical
subgraph.
\end{theorem}

The host is the classical graph commonly denoted $O_5$.
We do not claim a new construction of the host graph. Our argument
describes its 4-critical subgraphs and applies that description
to the question above.

\section{The graph and its criticality}

Use the vertex set $\{0,\ldots,8\}$, and write
\[
 L=\{0,1\},\qquad R=\{4,5\},\qquad
 Z=\{2,3,7\},\qquad u=6,\qquad v=8.
\]
The set $L\cup R$ induces $K_4$, and $Z$ induces $K_3$.
Both $u$ and $v$ are adjacent to every vertex of $Z$.
In addition, $u$ is adjacent to every vertex of $R$, and $v$
to every vertex of $L$. There are no other edges. In particular,
$uv$ is not an edge.

This is an Ore composition of two copies of $K_5$: delete the edge
$uv$ from one copy, split a vertex of the other copy into two
vertices with disjoint two-element neighborhoods, and identify
the split vertices with $u$ and $v$. The explicit description
above gives $6+3+6+4=19$ edges.

\begin{lemma}\label{lem:host}
The graph $G$ is noncomplete and 5-critical.
\end{lemma}

\begin{proof}
Suppose that $G$ has a proper four-coloring. The clique $L\cup R$
uses all four colors. The color of $u$ must occur on $L$, and
the color of $v$ must occur on $R$. Thus $u$ and $v$ have different
colors. Every vertex of the triangle $Z$ is adjacent to both,
leaving only two colors for $Z$, a contradiction.

A proper five-coloring, in vertex order $0,\ldots,8$, is
\[
 (0,1,0,1,2,3,4,2,3).
\]
It remains to check criticality. Permutations within $L$, within
$R$, and within $Z$ are automorphisms. Exchanging $L$ with $R$
and $u$ with $v$ is also an automorphism. These automorphisms
give five edge orbits, represented by the edges in
Table~\ref{tab:deletions}. Each row supplies a proper
four-coloring after deletion of that edge.

\begin{table}[ht]
\centering
\small
\begin{tabular}{cl}
\toprule
Deleted edge & Four-coloring in vertex order $0,\ldots,8$\\
\midrule
$01$ & $(0,0,3,2,3,2,1,0,1)$\\
$04$ & $(0,3,3,2,0,2,1,0,1)$\\
$23$ & $(0,2,0,0,3,1,2,3,1)$\\
$26$ & $(0,2,0,3,3,1,0,2,1)$\\
$46$ & $(0,3,3,2,1,2,1,0,1)$\\
\bottomrule
\end{tabular}
\caption{Four-colorings of representative edge deletions.}
\label{tab:deletions}
\end{table}

Every edge deletion is therefore four-colorable. Every vertex has
an incident edge, so every vertex deletion is four-colorable as
well, by restricting a coloring of an incident-edge deletion.
Every proper subgraph is contained in an edge or vertex deletion.
Finally, $u$ and $v$ are nonadjacent, so $G$ is not complete.
\end{proof}

\section{The 4-critical subgraphs}

We first record two elementary facts.

\begin{lemma}\label{lem:five}
Every graph on at most five vertices that is not three-colorable
contains $K_4$.
\end{lemma}

\begin{proof}
Suppose otherwise, and take a counterexample with the fewest
vertices. Every vertex has degree at least three: otherwise a
three-coloring after its deletion extends to the whole graph.
A counterexample on four vertices would be $K_4$, so there are
five vertices. Its complement has maximum degree at most one;
its missing edges therefore form a matching. If there is at
most one missing edge, the graph contains $K_4$. If there are two, give each
nonadjacent pair one color and the remaining vertex a third.
Both possibilities give a contradiction.
\end{proof}

\begin{lemma}\label{lem:spindles}
For every three-element subset $T$ of $L\cup R$ and every
two-element subset $P$ of $Z$, the induced graph
\[
 G[T\cup P\cup\{u,v\}]
\]
is a 4-critical Moser spindle.
\end{lemma}

\begin{proof}
By symmetry, write $T=\{a,b,c\}$ with $a,b\in L$ and $c\in R$,
and write $P=\{x,y\}$. The graph consists of two diamonds
(copies of $K_4$ with one edge deleted), sharing the vertex $v$,
together with one additional edge. The first diamond has
vertices $\{a,b,c,v\}$ and nonadjacent terminals $c,v$.
The second has vertices $\{x,y,u,v\}$ and nonadjacent terminals
$u,v$. The additional edge is $cu$. This is the usual
seven-vertex, eleven-edge description of the Moser spindle.

In every three-coloring of a diamond, its terminals have the
same color. The two diamonds would therefore force $c,v,u$
to have the same color, contradicting the edge $cu$.

For completeness, every edge deletion is three-colorable.
Deleting $cu$ permits both diamonds to be colored with all
three terminals equal. If an edge of a diamond is deleted,
that diamond admits a three-coloring with different terminal
colors. To check this last assertion, name the terminals $p,q$
and the other vertices $r,s$. After deleting $rs$, use
$(p,q,r,s)=(0,1,2,2)$. After deleting $pr$, use
$(p,q,r,s)=(0,1,0,2)$. Symmetry covers the other three edges.

Thus, after deleting an edge of either diamond, color its
terminals differently and color the other diamond with equal
terminals. Color permutations make the colorings agree at $v$;
the edge $cu$ then has different endpoint colors. A
three-coloring of an edge deletion also gives a four-coloring
of the original graph by assigning one endpoint a fourth color.
Every vertex has positive degree, so the edge-deletion checks
imply colorability of every proper subgraph.
\end{proof}

\begin{proof}[Proof of Theorem~\ref{thm:main}]
Lemma~\ref{lem:host} establishes the required properties of the
host. Its copies of $K_4$ are exactly those induced by
\[
 L\cup R,\qquad Z\cup\{u\},\qquad Z\cup\{v\}.
\]
There are no edges between $L\cup R$ and $Z$, and $u,v$ are
nonadjacent, so there are no other copies. Lemma~\ref{lem:spindles}
gives twelve further induced 4-critical subgraphs, since there
are $\binom43\binom32=12$ choices of $(T,P)$.
We show that this list is exhaustive.

Let $F$ be a 4-critical subgraph of $G$. If $F$ contains $K_4$,
criticality forces $F$ to equal that $K_4$. Assume henceforth
that $F$ contains no $K_4$.

For the coloring argument, add any missing host vertices as
isolated vertices. Split the edges of $F$ into a graph $X$ on
$(L\cup R)\cup\{u,v\}$ and a graph $Y$ on $Z\cup\{u,v\}$.
Their common vertex set consists of the independent terminals
$u,v$. By Lemma~\ref{lem:five}, $Y$ is three-colorable.

The graph $X$ has a three-coloring with $u$ and $v$ different.
Indeed, some edge $ab$ is missing from its four-vertex clique
side. Color $K_4-ab$ on $L\cup R$ by giving $a,b$ color $0$ and the
other two vertices colors $1,2$. If $a,b$ lie in different parts,
give $u$ the color of the remaining vertex of $L$ and $v$ the
color of the remaining vertex of $R$. If $a,b\in L$, give
$u$ color $0$ and $v$ color $1$; if $a,b\in R$, reverse these
terminal assignments. These colorings respect all permitted
terminal attachments and hence all edges of $X$.

Two three-colorings of $X$ and $Y$ can be joined after a color
permutation whenever they give the terminal pair the same
equality relation. Since $F$ is not three-colorable and $X$
permits different terminal colors, every three-coloring of $Y$
must give $u,v$ the same color. Consequently $Y+uv$ is not
three-colorable. Lemma~\ref{lem:five} gives a $K_4$ in $Y+uv$.
It must use the added edge, because $Y$ contains no $K_4$.
Thus $Y$ contains the diamond on $u,v$ and some pair
$P\subset Z$.

Since $Y$ is three-colorable with equal terminal colors, $X$
cannot admit such a coloring. Identifying $u$ and $v$ in $X$
therefore gives a non-three-colorable graph on five vertices.
By Lemma~\ref{lem:five}, it contains $K_4$. This $K_4$ must
contain the identified terminal, since the clique side of $X$
contains no $K_4$. It follows that $X$ contains a triangle on
some three-element subset $T\subset L\cup R$, together with
the permitted terminal attachment of each of its vertices.
The uniqueness of each vertex's permitted terminal neighbor
justifies lifting these three edges from the identified graph.

The triangle and these attachments, together with the diamond
in $Y$, form precisely $G[T\cup P\cup\{u,v\}]$.
Lemma~\ref{lem:spindles} shows that this is 4-critical.
Criticality of $F$ forces equality. The fifteen listed
subgraphs are distinct and induced, proving the theorem.
\end{proof}

\section{Formal verification}

The Lean source and reproducible build instructions are at
\begin{center}
\url{https://github.com/alexgocardinal/jt58-critical-subgraphs}.
\end{center}
The repository is private during manuscript preparation.

The formal theorem \texttt{JT58.jt58\_counterexample} proves
5-criticality with respect to arbitrary proper subgraphs,
noncompleteness, and the absence of a noninduced 4-critical
subgraph. The formalization uses Lean~4.35.0 and does not
require Mathlib. Explicit certificates check the edge and
vertex deletions of the host. For the remaining assertion,
a direct case proof covers all $2^{19}$ edge subsets using
fifteen induced non-three-colorable cores and 264
three-coloring certificates. Its 2,986 terminal branches
are checked by Lean's kernel.

Python scripts generate proof files, but their computations
are not trusted premises. The checked-in Lean files suffice
for verification. The axiom audit of the final theorem reports
exactly
\[
  \texttt{propext},\qquad
  \texttt{Classical.choice},\qquad
  \texttt{Quot.sound}.
\]
The proof uses neither a placeholder axiom nor a
native-computation axiom. Running \texttt{lake build}
checks the formal counterexample theorem. The exact
classification in Theorem~\ref{thm:main} is established
by the elementary argument above.

\begin{thebibliography}{9}

\bibitem{JT}
T.~R. Jensen and B.~Toft,
\emph{Graph Coloring Problems},
Wiley-Interscience, New York, 1995.
\href{https://doi.org/10.1002/9781118032497.ch5}{Chapter 5}.

\end{thebibliography}

\end{document}
